\section{Proof: Recovering Student Aleatoric Variance from a Noisy
Teacher}\label{proof-recovering-student-aleatoric-variance-from-a-noisy-teacher}

\subsection{1. Problem Setup}\label{problem-setup}

We have two predictors/sensors estimating the same underlying nominal
quantity \(y^*\):

\begin{itemize}
\tightlist
\item
  \textbf{Student:} observes \(x_1\)
\item
  \textbf{Teacher:} observes \(x_2\)
\end{itemize}

The student input \(x_1\) represents proprioceptive data. Its noise can
arise from sources such as slip and encoder noise.

The teacher input \(x_2\) represents visual features from the
FAST-LIO/ZED2 perception pipeline. Its noise can arise from the visual
features themselves, changing lighting conditions, and related
visual-sensing effects.

The student and teacher use two different sensor modalities. Their
inputs and noise sources are therefore assumed independent.

The student prediction is modeled as

\[
y_s = \mu_s(x_1) + \epsilon_s,
\]

where

\[
\mathbb{E}[\epsilon_s\mid x_1]=0,
\qquad
\operatorname{Var}(\epsilon_s\mid x_1)=\sigma_s^2(x_1).
\]

The teacher prediction is modeled as

\[
y_t = \mu_t(x_2) + \epsilon_t,
\]

where

\[
\mathbb{E}[\epsilon_t\mid x_2]=0,
\qquad
\operatorname{Var}(\epsilon_t\mid x_2)=\sigma_t^2(x_2).
\]

The teacher provides \(\sigma_t^2(x_2)\) for each training sample. The
student must learn \(\sigma_s^2(x_1)\).

\begin{center}\rule{0.5\linewidth}{0.5pt}\end{center}

\subsection{2. Same Nominal Output
Assumption}\label{same-nominal-output-assumption}

Assume that, for the corresponding physical state, the teacher and
student estimate the same nominal output:

\[
\mu_s(x_1)=\mu_t(x_2)=y^*.
\]

Therefore,

\[
y_s = y^*+\epsilon_s,
\]

\[
y_t = y^*+\epsilon_t.
\]

This assumption removes systematic mean disagreement between the teacher
and student. Their disagreement is therefore due only to their
respective noise terms.

\begin{center}\rule{0.5\linewidth}{0.5pt}\end{center}

\subsection{3. Define the Teacher--Student
Residual}\label{define-the-teacherstudent-residual}

Define

\[
r = y_t-y_s.
\]

Substituting the two observation models gives

\[
r
= (y^*+\epsilon_t)-(y^*+\epsilon_s),
\]

and hence

\[
\boxed{r=\epsilon_t-\epsilon_s.}
\]

The common nominal output cancels.

\begin{center}\rule{0.5\linewidth}{0.5pt}\end{center}

\subsection{4. Mean of the Residual}\label{mean-of-the-residual}

Because both noise terms are zero mean,

\[
\mathbb{E}[r\mid x_1,x_2]
=
\mathbb{E}[\epsilon_t-\epsilon_s\mid x_1,x_2].
\]

Under the assumed noise model,

\[
\boxed{\mathbb{E}[r\mid x_1,x_2]=0.}
\]

Thus the residual contains no systematic bias under the
same-nominal-output assumption.

\begin{center}\rule{0.5\linewidth}{0.5pt}\end{center}

\subsection{5. Variance of the Residual}\label{variance-of-the-residual}

From

\[
r=\epsilon_t-\epsilon_s,
\]

we obtain

\[
\operatorname{Var}(r\mid x_1,x_2)
=
\operatorname{Var}(\epsilon_t-\epsilon_s\mid x_1,x_2).
\]

In general,

\[
\operatorname{Var}(A-B)
=
\operatorname{Var}(A)
+
\operatorname{Var}(B)
-
2\operatorname{Cov}(A,B).
\]

Therefore,

\[
\operatorname{Var}(r\mid x_1,x_2)
=
\sigma_t^2(x_2)
+
\sigma_s^2(x_1)
-
2\operatorname{Cov}(\epsilon_t,\epsilon_s\mid x_1,x_2).
\]

Because the teacher and student have independent noise sources,

\[
\operatorname{Cov}(\epsilon_t,\epsilon_s\mid x_1,x_2)=0.
\]

Hence

\[
\boxed{
\operatorname{Var}(r\mid x_1,x_2)
=
\sigma_s^2(x_1)+\sigma_t^2(x_2).
}
\]

This is the central variance-addition result.

\begin{center}\rule{0.5\linewidth}{0.5pt}\end{center}

\subsection{6. Recovering the Student
Variance}\label{recovering-the-student-variance}

Rearranging the previous equation gives

\[
\boxed{
\sigma_s^2(x_1)
=
\operatorname{Var}(r\mid x_1,x_2)
-
\sigma_t^2(x_2).
}
\]

Therefore, conceptually,

\[
\boxed{
\text{student aleatoric variance}
=
\text{teacher–student residual variance}
-
\text{teacher aleatoric variance}.
}
\]

\begin{center}\rule{0.5\linewidth}{0.5pt}\end{center}

\subsection{7. The Training Likelihood}\label{the-training-likelihood}

Instead of explicitly estimating residual variance and subtracting
teacher variance, model the residual as

\[
r\mid x_1,x_2
\sim
\mathcal N\left(
0,
\sigma_s^2(x_1)+\sigma_t^2(x_2)
\right).
\]

The corresponding negative log-likelihood, ignoring constants
independent of the learned parameters, is

\[
\boxed{
\mathcal L
=
\frac{1}{2}
\left[
\frac{r^2}
{\sigma_s^2(x_1)+\sigma_t^2(x_2)}
+
\log\left(
\sigma_s^2(x_1)+\sigma_t^2(x_2)
\right)
\right].
}
\]

The teacher variance \(\sigma_t^2(x_2)\) is supplied for each sample.
The network predicts only \(\sigma_s^2(x_1)\).

\begin{center}\rule{0.5\linewidth}{0.5pt}\end{center}

\subsection{8. What Does the Loss
Learn?}\label{what-does-the-loss-learn}

For a fixed student input \(x_1\), define

\[
v=\sigma_s^2(x_1)
\]

and

\[
t=\sigma_t^2(x_2).
\]

The expected population loss is

\[
J(v\mid x_1)
=
\mathbb{E}
\left[
\frac{1}{2}
\left(
\frac{r^2}{v+t}
+
\log(v+t)
\right)
\middle|x_1
\right].
\]

Differentiate with respect to \(v\):

\[
\frac{\partial J}{\partial v}
=
\frac{1}{2}
\mathbb{E}
\left[
-\frac{r^2}{(v+t)^2}
+\frac{1}{v+t}
\middle|x_1
\right].
\]

Combining terms,

\[
\boxed{
\frac{\partial J}{\partial v}
=
\frac{1}{2}
\mathbb{E}
\left[
\frac{v+t-r^2}{(v+t)^2}
\middle|x_1
\right].
}
\]

At the population optimum,

\[
\mathbb{E}
\left[
\frac{v+t-r^2}{(v+t)^2}
\middle|x_1
\right]=0.
\]

\begin{center}\rule{0.5\linewidth}{0.5pt}\end{center}

\subsection{9. Show That the True Student Variance Is an
Optimum}\label{show-that-the-true-student-variance-is-an-optimum}

Let the true student aleatoric variance be

\[
v^*=\sigma_{s,\mathrm{true}}^2(x_1).
\]

From the generative model,

\[
\mathbb{E}[r^2\mid x_1,x_2]
=
v^*+t.
\]

Evaluate the expected gradient at \(v=v^*\):

\[
\left.\frac{\partial J}{\partial v}\right|_{v=v^*}
=
\frac12
\mathbb{E}
\left[
\frac{v^*+t-r^2}{(v^*+t)^2}
\middle|x_1
\right].
\]

By the law of iterated expectations, we may condition first on \(x_2\):

\[
\mathbb{E}
\left[
\frac{v^*+t-r^2}{(v^*+t)^2}
\middle|x_1
\right]
=
\mathbb{E}_{x_2\mid x_1}
\left[
\mathbb{E}
\left[
\frac{v^*+t-r^2}{(v^*+t)^2}
\middle|x_1,x_2
\right]
\right].
\]

For fixed \((x_1,x_2)\), the denominator is fixed, so

\[
\mathbb{E}
\left[
\frac{v^*+t-r^2}{(v^*+t)^2}
\middle|x_1,x_2
\right]
=
\frac{
v^*+t-\mathbb{E}[r^2\mid x_1,x_2]
}{(v^*+t)^2}.
\]

Using

\[
\mathbb{E}[r^2\mid x_1,x_2]=v^*+t,
\]

we obtain

\[
\frac{v^*+t-(v^*+t)}{(v^*+t)^2}=0.
\]

Therefore,

\[
\boxed{
\left.\frac{\partial J}{\partial v}\right|_{v=v^*}=0.
}
\]

Thus the true student aleatoric variance is a stationary population
solution of the correctly specified Gaussian likelihood.

Under the usual identifiability, sufficient-data, optimization, and
model-capacity assumptions, maximum-likelihood training is therefore
consistent with learning

\[
\boxed{
\sigma_s^2(x_1)\rightarrow\sigma_{s,\mathrm{true}}^2(x_1).
}
\]

\begin{center}\rule{0.5\linewidth}{0.5pt}\end{center}

\subsection{10. What Happens to the Varying Teacher
Variance?}\label{what-happens-to-the-varying-teacher-variance}

The teacher variance may vary strongly with \(x_2\):

\[
\sigma_t^2=0.001,\;0.1,\;0.03,\;1.2,\ldots
\]

The student does \textbf{not} need to predict \(x_2\), learn
\(\sigma_t^2(x_2)\), or manually subtract an average teacher variance
before training.

For every sample, the supplied teacher variance appears directly in

\[
\sigma_s^2(x_1)+\sigma_t^2(x_2).
\]

Consequently, each residual is interpreted relative to the amount of
teacher noise present in that particular sample.

Across the training distribution, optimization performs the required
averaging/marginalization implicitly.

\begin{center}\rule{0.5\linewidth}{0.5pt}\end{center}

\subsection{11. Relation to the Marginal Variance
Equation}\label{relation-to-the-marginal-variance-equation}

If one additionally assumes

\[
x_1\perp x_2,
\]

then

\[
\mathbb{E}[\sigma_t^2(x_2)\mid x_1]
=
\mathbb{E}[\sigma_t^2(x_2)].
\]

Marginalizing the residual second moment over \(x_2\) gives

\[
\mathbb{E}[r^2\mid x_1]
=
\sigma_s^2(x_1)
+
\mathbb{E}[\sigma_t^2(x_2)].
\]

Hence

\[
\boxed{
\sigma_s^2(x_1)
=
\mathbb{E}[r^2\mid x_1]
-
\mathbb{E}[\sigma_t^2(x_2)].
}
\]

This equation describes the \textbf{marginal moment relationship}. It is
not necessary to explicitly compute this subtraction when training with
the sample-wise Gaussian likelihood above.

\begin{center}\rule{0.5\linewidth}{0.5pt}\end{center}

\subsection{12. Assumptions Required for the
Proof}\label{assumptions-required-for-the-proof}

The result depends on the following assumptions:

\begin{enumerate}
\def\labelenumi{\arabic{enumi}.}
\item
  \textbf{Same nominal output}

  \[
  \mu_s(x_1)=\mu_t(x_2)=y^*.
  \]
\item
  \textbf{Zero-mean student noise}

  \[
  \mathbb{E}[\epsilon_s\mid x_1]=0.
  \]
\item
  \textbf{Zero-mean teacher noise}

  \[
  \mathbb{E}[\epsilon_t\mid x_2]=0.
  \]
\item
  \textbf{Independent student and teacher noise sources}

  Because \(x_1\) and \(x_2\) come from two different sensor modalities
  (proprioceptive and visual), their noise sources are assumed
  independent.

  \[
  \operatorname{Cov}(\epsilon_s,\epsilon_t\mid x_1,x_2)=0.
  \]
\item
  \textbf{Teacher covariance is calibrated and known for each training
  sample}

  \[
  \sigma_t^2(x_2)
  =
  \operatorname{Var}(\epsilon_t\mid x_2).
  \]
\item
  \textbf{The likelihood model is correctly specified}, or sufficiently
  close to the actual residual distribution for the intended
  interpretation.
\item
  \textbf{Independent student and teacher inputs for the additional
  global-average simplification in Section 12 only}

  Because \(x_1\) and \(x_2\) arise from different sensor modalities,
  assume

  \[
  x_1\perp x_2.
  \]
\end{enumerate}

The independence of \(x_1\) and \(x_2\) is \textbf{not required} for the
basic variance-addition equation or for the sample-wise likelihood
argument. Independence of the \textbf{noise terms} is the key
requirement for variance addition.

\begin{center}\rule{0.5\linewidth}{0.5pt}\end{center}

\subsection{13. Final Result}\label{final-result}

Under these assumptions,

\[
\boxed{
\operatorname{Var}(y_t-y_s\mid x_1,x_2)
=
\sigma_s^2(x_1)+\sigma_t^2(x_2)
}
\]

and therefore

\[
\boxed{
\sigma_s^2(x_1)
=
\operatorname{Var}(y_t-y_s\mid x_1,x_2)
-
\sigma_t^2(x_2).
}
\]

Training with

\[
\boxed{
\mathcal L
=
\frac12
\left[
\frac{(y_t-y_s)^2}
{\sigma_s^2(x_1)+\sigma_t^2(x_2)}
+
\log\left(\sigma_s^2(x_1)+\sigma_t^2(x_2)\right)
\right]
}
\]

allows the student to recover its own aleatoric variance without
explicitly estimating or subtracting the teacher's average variance,
provided the model assumptions hold.